\section{Proofs for world valuation and residual substitution}
\label{app:stage23}
\subsection{World-number optimization}
\label{app:world-proof}
\begin{proof}[Proof of \cref{lem:world-curvature}]
Eliminate $W$ to obtain $C_{\rm tot}^{\theta}c^{\alpha-\theta}$ on
$[u,C_{\rm tot}]$. If $\alpha>\theta$, its maximum is attained at
$c=C_{\rm tot}$ and $W=1$. If $\alpha<\theta$, it is attained at
$c=u$ and $W=C_{\rm tot}/u$. Equality of the exponents makes every
split tie. This proves \cref{eq:world-value}, including $C_{\rm tot}=u$.
In either branch $V(u)=u^\alpha$ and
$V(u(1+C))/V(u)=(1+C)^v$. Since $v>0$, the stated affine
normalization is $[(1+C)^v-1]/v=F_{1-v}(C)$.
An affine rescaling of the computational component requires the residual
weights to be expressed in the same utility units; it is not neutral if
those weights are held in unrelated units.

For logarithmic within-world welfare put $L=\ln(C_{\rm tot}/u)$ and
$v_c=\ln(c/u)$. The objective is $\exp[\theta(L-v_c)]v_c$ on $[0,L]$.
At $\theta=0$ it increases to $L$. At $\theta>0$ its derivative has the
sign of $1-\theta v_c$, so $v_c^*=\min(L,1/\theta)$. For
$L\geq1/\theta$ substitution gives the claimed power expression;
below that boundary the maximum is $L$ with one world. Consequently
$\theta=1$ gives linear value only above $\exp(1)u$ in this logarithmic
case. No $\eta=1$ value follows from the positive-$\alpha$ power model.

With integer $W$, eliminate $c$ instead: the objective is
$C_{\rm tot}^\alpha W^{\theta-\alpha}$ for
$1\leq W\leq\lfloor C_{\rm tot}/u\rfloor$. The same endpoint argument
gives the stated integer formula when $\theta>\alpha$, and $W=1$
when $\theta<\alpha$. Since $\lfloor t\rfloor/t\to1$, the asymptotic
exponent is unchanged. For regular counts $W_d$, the largest admissible
count is between fixed positive multiples of $C_{\rm tot}/u$ for all
sufficiently large $C_{\rm tot}$, because $W_{d+1}=bW_d+1$.
These bounds give the same logarithmic growth exponent; they do not
give a smooth exact power function or a ratio converging to one.
\end{proof}

\subsection{The optimized no-replication face}
\label{app:replication-proof}
\begin{proof}[Proof of \cref{prop:no-replication}]
The feasible set is a compact simplex; the objective is continuous and
concave, so a face optimizer exists. Differentiability and the positive
budget give the usual multiplier conditions, including a multiplier
$\lambda$ for the budget and complementary multipliers for coordinates
at zero. Equivalently, for each face coordinate $v$ of unit resource
price $p_v$, its marginal utility is at most $p_v\lambda$, with equality
at positive coordinates. This can also be obtained directly by exchanges
between any two positive-resource uses. At least one use is positive.
The local derivative is positive, so the resulting $\lambda$ is positive.

At a point with $x=y=0$, the two missing marginal utilities are precisely
the left sides of \cref{eq:no-replication}. If these are at most their
prices times the same $\lambda$, then for every feasible point $v$ the
supporting-hyperplane bound is
\[
 U(v)-U(v^*)\leq\nabla U(v^*)\cdot(v-v^*)\leq0.
\]
The second inequality follows from the exhausted budget, equality in
positive face coordinates, and the non-positive reduced derivatives at
zero coordinates. Strict reduced derivatives for $x,y$ make this bound
strict whenever either is positive. Conversely, a full optimum satisfies
the full multiplier conditions; restricting that multiplier to the face
gives the stated existential condition. Thus the statement remains
correct when several restricted optima or multipliers are possible.
For $z^*>0$, complementary slackness fixes
$\lambda=F'(C^*)a\ell'(z^*)$. If $z^*=0$, only
$F'(C^*)a\ell'(0)\leq\lambda$ is required, assuming the derivative is
finite. The potentially non-differentiable endpoint is outside the
present version of the proposition and can instead be treated with
supergradients. If backups are unavailable, omit their coordinate and
retain $R'(0)$ throughout.
\end{proof}

\begin{proof}[Proof of \cref{cor:bundle} and the displayed reductions]
For $w^*>0$, $\epsilon_II'(w^*)=p_m\lambda$; subtracting it from
\cref{eq:no-replication} proves \cref{eq:observation-net} when no
backup exists. For $b^*>0$, subtract $p_b\lambda$ in the same manner.
With both channels positive the industrial reduced derivative is
$F'(C^*)r_h+(p_m+p_b-p_h)\lambda$, strictly positive when
$p_h<p_m+p_b$. At equality and $r_h=0$, non-decreasing concavity and
$G_h'(0)=0$ imply $G_h(y)=0$ for $y\geq0$. For
$0\leq\epsilon\leq\min(w^*,b^*)$, replace
$(y,w,b_{\rm bk})=(0,w^*,b^*)$ by
$(\epsilon,w^*-\epsilon,b^*-\epsilon)$. The budget, $C$, $x+y+b_{\rm bk}$,
and $x+y+w$ are unchanged, proving the tie statement. With a zero
observation channel only $\epsilon_II'(0)\leq p_m\lambda$ is known;
division by positive $\lambda$ defines $c_I\leq p_m$ and gives its exact
shadow service cost. The backup argument is identical.
Finally, at $q=1$ the experiential marginal vanishes identically.
Substitution proves \cref{eq:replication-reduced}; its sign and tie cases
follow from non-negative returns and positive $F'(C^*),\lambda$.
\end{proof}

\subsection{Expected value at risk}
\label{app:risk-proof}
\begin{proof}[Proof of \cref{prop:endogenous-risk}]
Let $k(x)=x/C(x)^\gamma$. Since $C(0)=\Omega>0$,
$k'(0)=\Omega^{-\gamma}$: the term differentiating the denominator is
multiplied by $x=0$. Also $C'(0)=r-pa\ell'(B)$ and
$\varphi(0)=0$. The complete product derivative is
\[
 U_R'(0)=(1-\vartheta)F_\eta'(\Omega)[r-pa\ell'(B)]
   +\vartheta\varphi'(0)F_\eta(\Omega)\Omega^{-\gamma}.
\]
A positive derivative improves along a feasible sufficiently small
positive $x$, so an inward optimum requires this expression to be
non-positive. Divide by $(1-\vartheta)F_\eta'(\Omega)>0$, rearrange,
and use $a\ell'(B)=\Omega/\Gamma$ to obtain \cref{eq:risk-floor}.
For $\eta\ne1$ exact algebra yields
\[
 K_\eta(\Omega)=
 \frac{(1+\Omega)-(1+\Omega)^\eta}{(1-\eta)\Omega}.
\]
For $0\leq\eta<1$, $(1+\Omega)/\Omega\to1$ and
$(1+\Omega)^\eta/\Omega\to0$, giving the claimed limit.
In particular, for every $\varepsilon>0$ there is $\Omega_\varepsilon$
such that
$|K_\eta(\Omega)-(1-\eta)^{-1}|<\varepsilon$ for all
$\Omega\geq\Omega_\varepsilon$. This supplies a rigorous two-sided
asymptotic replacement even when $\Gamma$ or the hazard varies; it is
not an exact finite-output inequality with the limiting constant.
The logarithmic expression follows directly from its definition, and
its ratio to $\ln\Omega$ tends to one. For completeness, when $\eta>1$,
$K_\eta(\Omega)\sim\Omega^{\eta-1}/(\eta-1)$.

For seeds, a positive limiting $K_\eta$ and bounded $p$ force hazard
odds times $\Gamma$ to stay bounded. At a fixed budget, with $\Omega\to
\infty$ and fixed $r$, the remote term tends to zero; taking limits
and solving the resulting odds inequality proves \cref{eq:hazard-ceiling}.
For mirrors at bounded positive budgets,
$\Gamma r/\Omega\to0$ and
$\Gamma K_\eta(\Omega)/\Omega\to0$ for $\eta\leq1$.
With logarithmic value and seeds, the same exact inequality instead
bounds hazard odds times $\Gamma\ln\Omega$. These consequences are
necessary local comparisons unless additional global structure is
supplied. At $\vartheta=0$ the risk term vanishes. At
$\vartheta=1$, division by surviving marginal value is invalid; with
$\varphi'(0)>0$ and $F_\eta(\Omega)>0$ the original derivative is
positive and the inward corner fails. If $\varphi'(0)=0$, this last
first-order conclusion need not hold. The scaling specification itself
requires $\Omega>0$; no $0/0$ backup is defined at zero computation.
\end{proof}

\subsection{Proof of the growing-budget separation}
\label{app:growing-norep-proof}
The restricted objective is strictly concave in $(z,w,b)$ because the
local log-composition and the two positive exponential utilities are
strictly concave. It has a unique budget-exhausting optimum. If $z=0$,
the local marginal utility is $B^2/d_0$, whereas either positive service
coordinate has marginal utility per resource at most a fixed constant.
The budget multiplier conditions therefore rule out $z=0$ for large $B$.
With $z>0$,
\[
 \lambda=\frac{B^2}{(1+z)[d_0+B^2\ln(1+z)]}.
\]
We claim $z\to\infty$. Otherwise along a subsequence $z$ is bounded,
so this expression is bounded below by a positive constant. The service
optimality conditions then imply
\[
 w=s_I\max\{0,\ln[J_I/(s_Ip_m\lambda)]\},\qquad
 b=s_R\max\{0,\ln[J_R/(s_Rp_b\lambda)]\},
\]
which bound both services on that subsequence. This contradicts
$B=z+p_mw+p_bb\to\infty$. Hence $z\to\infty$ and
$\lambda\sim[(1+z)\ln(1+z)]^{-1}\to0$. Both services are eventually
positive, and their formulas give $w+b=O(\ln z)$ and
$B=z+O(\ln z)$. Consequently $z/B\to1$.
With $q=1$, \cref{eq:no-replication} reduces to the net-price conditions.
Since $F'(C^*)/\lambda=(1+z)/B^2$, each required computational premium
$r_j(1+z)/B^2$ tends to zero. The fixed positive margins
$p_j-p_m-p_b$ make both conditions strict eventually, excluding replication
in every optimizer. Finally, at the all-inward corner the service slopes
$J_I/s_I$ and $J_R/s_R$ are fixed and positive, and bounded project prices
contradict the growing logarithmic residual floor. This proves the
claimed separation without assuming any resource-acquisition dynamics.


\subsection{The no-replication construction under a finite ceiling}
\label{app:norep-ceiling-proof}
Keep the fixed parameters, logarithmic production and value, exponential
service utilities, $q=1$, and positive net prices of
\cref{cor:growing-no-replication}, but allow any positive productivity
path with $a\leq a_{\max}<\infty$. The restricted face remains strictly
concave and budget-exhausting. At least one of the expenditures
$z,p_mw,p_bb$ is at least $B/3$. If it is $z$, its positive-coordinate
multiplier condition yields
\[
 \lambda=\frac{a}{(1+z)[d_0+a\ln(1+z)]}
 \leq\frac{3a_{\max}}{d_0B}.
\]
If it is $p_mw$, put $u=w/s_I$ and use $u e^{-u}\leq e^{-1}$:
$\lambda p_mw=J_Iu e^{-u}\leq J_I/e$, whence
$\lambda\leq3J_I/(eB)$. The backup case gives $3J_R/(eB)$.
Consequently, with the fixed finite constant
$K=3\max\{a_{\max}/d_0,J_I/e,J_R/e\}$,
\begin{equation}
 0<\lambda\leq K/B,\qquad
 \frac{F'(C^*)}{\lambda}
 \geq\frac{B}{K[d_0+a_{\max}\ln(1+B)]}\longrightarrow\infty.
 \label{eq:norep-ceiling-bound}
\end{equation}
This uniform proof includes $z=0$ and needs no positive lower bound on
$a$. Since $\lambda\to0$, each service's strictly positive marginal
value at zero forces $w,b>0$ eventually. The reduced test is then
$r_jF'(C^*)/\lambda\leq\Delta_j$, where
$\Delta_j=p_j-p_m-p_b>0$ is fixed. Every $r_j>0$ eventually fails it.
No no-replication face optimum is then a full optimum.

If $z>0$, $\lambda=F'(C^*)a\ell'(z)$ gives the exact identity
$F'(C^*)/\lambda=1/[a\ell'(z)]=(1+z)/a$ and the exact condition
$a\ell'(z)\geq r_j/\Delta_j$ at the optimized face point.
The ceiling provides the necessary inequality with $a_{\max}$, not a
sufficient condition for an arbitrary actual productivity or an
unoptimized choice of $z$. If $r_j=0$, that channel's reduced inequality
is strict for $\Delta_j>0$; if both replication slopes vanish, both
channels pass. Non-positive net prices retain the separate sign and tie
cases in \cref{cor:bundle}. At $z=0$, use the multiplier test rather than
the interior identity. For the pure-computation comparison in
\cref{eq:budgetcrit}, the same logarithmic marginal production
$a/(1+z)$ declines with growing local investment; the present calculation
prices the optimized substitute services as well.
